\section{Reference Charge-Pump PFC Rectifier}
\label{sec:bg-ref2020}

The analytical method was developed by Ammar et al. for a charge-pump resonant AC--DC converter used as an offline PFC rectifier \cite{ammar2020chargepump}. Adding a capacitor and diode to the class-DE circuit shapes the input current without a separate current-loop PFC controller. Their analysis connects capacitor charge transfer with resonant gain and DC-link energy storage.

The reference converter was operated near \SI{1}{\mega\hertz}, with output power up to \SI{50}{\watt}. The reported values were PF of 0.99, input-current THD of 8.6\%, and efficiency up to 88\%. These results demonstrate the principle and provide a starting point for the design calculations.

\begin{table}[htbp]
\centering
\caption{Reported performance of the 2020 reference rectifier \cite{ammar2020chargepump}.}
\label{tab:bg-ref2020-benchmark}
\begin{tabularx}{\textwidth}{@{}L{0.42\textwidth}Y@{}}
\toprule
Quantity & Reported result\\
\midrule
Function & Resonant AC--DC PFC rectifier\\
Switching frequency & Approximately \SI{1}{\mega\hertz}\\
Maximum output power & \SI{50}{\watt}\\
Power factor & 0.99\\
Input-current THD & 8.6\%\\
Peak efficiency & 88\%\\
\bottomrule
\end{tabularx}
\end{table}

Frequency changes affect both charge transfer and resonant gain. The prototype uses different magnetics, isolation, a downstream regulator, and sampled burst control, so the reference efficiency cannot be assumed for this supply.

The charge-transfer analysis is retained. The following isolated-driver architecture provides the closer precedent, with separate pump and output paths sharing one half bridge.
